\documentclass[10pt]{article}

% Draft for Digital Discovery (RSC). Generic article class for portability;
% swap in the RSC template at submission time.
%
% WORKING RULES (audit 2026-07-14):
% - Every number in prose comes from a macro in numbers.tex (generated by
%   paper/make_tables.py from committed results JSONs). Never type digits.
% - Significance claims must exist in a committed results file, tagged with
%   metric and pair, Holm-corrected within their family.
% - Comparisons only on identical support sets.

\usepackage[a4paper,margin=2.5cm]{geometry}
\usepackage{graphicx}
\usepackage{booktabs}
\usepackage{siunitx}
\usepackage{amsmath}
\usepackage[numbers,sort&compress]{natbib}
\usepackage[colorlinks=true,linkcolor=blue,citecolor=blue,urlcolor=blue]{hyperref}
\usepackage{xcolor}

% AUTO-GENERATED by paper/make_tables.py — do not edit by hand.
% Every number quoted in prose must come from a macro in this file.
\newcommand{\nRetrievalQueries}{400}
\newcommand{\corpusSize}{7,498}
\newcommand{\jaccFormulaSro}{0.362}
\newcommand{\magpieFormulaSro}{0.162}
\newcommand{\stoichFormulaSro}{0.324}
\newcommand{\randomFormulaSro}{0.017}
\newcommand{\hybridFormulaSro}{0.268}
\newcommand{\oracleFormulaSro}{0.595}
\newcommand{\jaccMeanRegret}{0.234}
\newcommand{\jaccFracLargeRegret}{52\%}
\newcommand{\jaccSolidStateFormulaSro}{0.492}
\newcommand{\oracleSolidStateFormulaSro}{0.654}
\newcommand{\gapSolidState}{0.162}
\newcommand{\jaccSolGelFormulaSro}{0.232}
\newcommand{\oracleSolGelFormulaSro}{0.537}
\newcommand{\gapSolGel}{0.305}
\newcommand{\jaccElementSro}{0.534}
\newcommand{\magpieElementSro}{0.325}
\newcommand{\nPlanningCases}{100}
\newcommand{\skyFOne}{0.497}
\newcommand{\skyFOneCI}{[0.411, 0.583]}
\newcommand{\skyElemJacc}{0.813}
\newcommand{\skyMethodAcc}{0.53}
\newcommand{\skyFailures}{2}
\newcommand{\llmFOne}{0.433}
\newcommand{\llmFOneCI}{[0.352, 0.517]}
\newcommand{\llmElemJacc}{0.811}
\newcommand{\llmMethodAcc}{0.51}
\newcommand{\llmFailures}{0}
\newcommand{\ruleFOne}{0.444}
\newcommand{\ruleFOneCI}{[0.360, 0.537]}
\newcommand{\ruleElemJacc}{0.821}
\newcommand{\ruleMethodAcc}{0.50}
\newcommand{\ruleFailures}{0}
\newcommand{\retrJaccFOne}{0.448}
\newcommand{\retrJaccFOneCI}{[0.374, 0.522]}
\newcommand{\retrJaccElemJacc}{0.754}
\newcommand{\retrJaccMethodAcc}{0.59}
\newcommand{\retrJaccFailures}{0}
\newcommand{\retrMagpieFOne}{0.233}
\newcommand{\retrMagpieFOneCI}{[0.174, 0.299]}
\newcommand{\retrMagpieElemJacc}{0.370}
\newcommand{\retrMagpieMethodAcc}{0.29}
\newcommand{\retrMagpieFailures}{42}
\newcommand{\skyVsRetrJaccFOnePRaw}{0.213}
\newcommand{\skyVsRetrJaccFOnePHolm}{0.852}
\newcommand{\skyVsRetrJaccFOneEffect}{-0.12}
\newcommand{\skyVsRuleFOnePRaw}{0.096}
\newcommand{\skyVsRuleFOnePHolm}{0.482}
\newcommand{\skyVsRuleFOneEffect}{-0.12}
\newcommand{\skyVsLlmFOnePRaw}{0.011}
\newcommand{\skyVsLlmFOnePHolm}{0.067}
\newcommand{\skyVsLlmFOneEffect}{-0.15}
\newcommand{\tempCommonN}{41}
\newcommand{\ruleTempMaeCommon}{317}
\newcommand{\skyTempMaeCommon}{383}
\newcommand{\llmTempMaeCommon}{340}


\newcommand{\todo}[1]{\textcolor{red}{[TODO: #1]}}

\title{Recipe transferability is the unmeasured bottleneck of
retrieval-grounded inorganic synthesis planning}

\author{%
  Kinga Mastej$^{1}$ \and Hyunsoo Park$^{1}$ \\[4pt]
  \normalsize $^{1}$Department of Materials, Imperial College London, UK \\
  \normalsize \todo{confirm author order, affiliations, and any additional authors}%
}
\date{\today\ --- draft, not for distribution}

\begin{document}
\maketitle

\begin{abstract}
Retrieval sits at the core of modern synthesis-planning systems --- learned
precursor recommenders, retrieval-augmented language models, and LLM agents
alike ground their predictions in ``similar'' known materials --- yet no
published evaluation measures whether the retrieved materials actually have
transferable recipes. We introduce a transferability benchmark over a
multi-route text-mined corpus of \corpusSize\ inorganic synthesis recipes:
retrieval is scored by formula-level recipe overlap (a non-circular metric
over normalised precursor sets), calibrated against both a trivial
element-overlap baseline and the oracle ceiling computable by exhaustive
search. We find (i)~learned composition embeddings (Magpie) score
\magpieFormulaSro\ formula-SRO@5, well below raw element-set Jaccard
(\jaccFormulaSro); (ii)~even the best retriever leaves most of the
achievable headroom on the table (oracle \oracleFormulaSro; mean regret
\jaccMeanRegret, with \jaccFracLargeRegret\ of queries losing $>$0.2);
and (iii)~the gap concentrates in sol-gel synthesis
(gap \gapSolGel\ vs \gapSolidState\ for solid-state), pointing to
route-conditioned precursor-form selection as the learnable target. In a
downstream planning benchmark on \nPlanningCases\ leakage-guarded held-out
targets, a tool-grounded agent (SKY, Synthesis Knowledge Yield) attains the
highest precursor F1 point estimate (\skyFOne\ vs \retrJaccFOne\ for
retrieval-copy), but the difference is not significant after Holm
correction ($p_{\mathrm{Holm}}=\skyVsRetrJaccFOnePHolm$) --- consistent
with recent evidence that retrieval quality, not agentic machinery, is the
operative bottleneck. We release the benchmark, oracle tooling, frozen test
sets, and all cached model outputs.
\end{abstract}

\section{Introduction}

Predicting how to synthesise a target inorganic material remains a
bottleneck of computational materials discovery: high-throughput screening
proposes candidates far faster than laboratories can devise recipes for
them~\cite{jain2013materials,szymanski2023alab}. Nearly every data-driven
approach to this problem --- learned precursor
recommenders~\cite{he2023precursor,noh2024retrievalretro},
retrieval-augmented language models~\cite{iclr2026benchmark}, and LLM
agents~\cite{bran2024chemcrow,boiko2023coscientist} --- shares one
architectural commitment: \emph{retrieve} materials similar to the target,
then transfer or adapt their known recipes.

The retrieval step itself, however, is essentially unexamined. Retrieval
quality is usually reported through downstream proxy metrics (top-$k$
precursor accuracy) on solid-state-only corpora, which cannot answer three
basic questions. Do retrieved ``similar'' materials actually have
transferable recipes? How far is any given retriever from the best
achievable on its corpus? And where does retrieval fail --- uniformly, or
in identifiable regimes? Recent benchmarking already hints the stakes are
high: retrieval augmentation is the only intervention that consistently
improves LLM synthesis prediction, while added agentic machinery often
does not~\cite{iclr2026benchmark}.

We answer these questions directly. Our contributions are:
\begin{enumerate}
\item \textbf{A transferability benchmark} over a multi-route corpus of
  \corpusSize\ text-mined recipes~\cite{kononova2019text}: retrieval is
  scored by \emph{formula-level synthesis recipe overlap} (formula-SRO) ---
  Jaccard similarity of normalised precursor formula sets --- a metric that
  is not circular for element-overlap retrievers, unlike element-level SRO.
\item \textbf{Oracle calibration.} For every query we compute the ceiling
  achievable by exhaustive search, turning retriever scores into
  \emph{regret}. Every table reports the trivial-baseline floor and the
  oracle ceiling.
\item \textbf{Failure analysis.} The oracle gap is not uniform: it
  concentrates in sol-gel chemistry, and (in a route-oracle ablation)
  conditioning on the true route closes only a small fraction of it ---
  the remainder is precursor-\emph{form} selection within a route
  \todo{reimplement route-oracle ablation in src/evaluation before
  quoting its numbers}.
\item \textbf{A leakage-guarded downstream planning benchmark} on
  \nPlanningCases\ held-out targets comparing rule, retrieval-copy,
  ungrounded-LLM, and agentic prediction under identical support sets and
  Holm-corrected statistics.
\end{enumerate}

SKY (Synthesis Knowledge Yield), the LLM agent whose development motivated
this study, appears here as an \emph{application} of the benchmark rather
than its subject: its measured behaviour illustrates that better agents do
not compensate for weak retrieval.

\section{Related work}

\paragraph{Learned precursor recommenders.} He \emph{et
al.}~\cite{he2023precursor} learn materials similarity for precursor
recommendation (solid-state, 29{,}900 recipes, top-5 success $\ge$82\%);
Retrieval-Retro~\cite{noh2024retrievalretro} combines learned and
thermodynamic retrievers over 33k recipes with year-based splits. Neither
measures transferability of the retrieved set on a non-circular metric,
reports oracle ceilings, calibrates against trivial baselines, nor covers
multi-route corpora. We evaluate on a multi-route corpus and will run both
systems inside our harness in follow-up work (\todo{T3 baselines rerun
before submission}).

\paragraph{LLM-based synthesis prediction.} Fine-tuned and off-the-shelf
LLMs predict precursors and conditions~\cite{prein2025ranking,msp_llm2026};
Prein \emph{et al.}~\cite{prein2025ranking} report top-1 precursor accuracy
of 53.8\% and calcination/sintering temperature MAE below
\SI{126}{\celsius} under a per-operation, solid-state protocol. Our
temperature figures are not comparable to theirs (we score max-heating
across mixed routes with regex-extracted ground truth); reconciling the
protocols is planned work, and we make no cross-paper temperature claims.

\paragraph{Agents and autonomous laboratories.}
ChemCrow~\cite{bran2024chemcrow}, Coscientist~\cite{boiko2023coscientist},
and A-Lab~\cite{szymanski2023alab} are demonstration-evaluated rather than
benchmarked against baselines. The ICLR 2026 augmentation
study~\cite{iclr2026benchmark} finds retrieval augmentation the only
consistently helpful strategy for synthesis prediction; our downstream
results independently support that conclusion. SKY originated at the 2025
LLM Hackathon for Applications in Materials Science and
Chemistry~\cite{hackathon2025outcomes}.

\section{Benchmark design}
\label{sec:benchmarks}

\subsection{Transferability benchmark}
For \nRetrievalQueries\ stratified query materials over the
\corpusSize-material corpus, a retriever returns $k$ neighbours;
the score is formula-SRO@$k$: the mean Jaccard similarity between the
query's and each neighbour's \emph{normalised precursor formula sets}
(pymatgen reduced formulas; hydrate forms distinct from anhydrous). We
also report element-level SRO for continuity with prior work, method
consistency rate (MCR), and per-query wall-clock latency. The \emph{oracle
ceiling} is the mean over the $k$ highest formula-set Jaccards achievable
for each query by exhaustive leave-one-out search; \emph{regret} is oracle
minus achieved. Retrievers: Magpie composition
embeddings~\cite{ward2016magpie} ($k$-NN in standardised feature space),
random floor, element-set Jaccard, stoichiometric-vector cosine, formula
TF--IDF, and a reciprocal-rank-fusion hybrid.

\subsection{Planning benchmark}
\nPlanningCases\ held-out targets, stratified 50/50 across the two route
families present in the corpus (solid-state, sol-gel), each with parseable
precursor ground truth. The stratification balances a corpus prior of
roughly 86\% solid-state; we therefore also interpret method-prediction
results against that prior (Section~\ref{sec:results}). \textbf{The test
set is a committed file, not a seed}: the seeded builder's candidate pool
depends on the installed pymatgen version's formula parsing, and an
environment change reshuffles the sample (we verified only 17/100 targets
are recovered under a newer pymatgen). \textbf{Leakage guard:} held-out
formulas are removed from every retrieval corpus and recipe store.

Methods predict precursors, route, and maximum heating temperature; we
score precursor element Jaccard, precursor formula F1, method accuracy,
and temperature absolute error. Aggregates carry bootstrap 95\% CIs;
pairwise comparisons use Wilcoxon signed-rank tests with Holm correction
within each metric family; temperature MAE is compared only on the common
subset of cases scoreable for every method ($n=\tempCommonN$), with
per-method coverage disclosed. Compared methods: a fitted rule baseline,
retrieval-copy (element-Jaccard and Magpie variants), an ungrounded LLM,
and the SKY agent (tool loop over similarity search and leakage-aware
recipe lookup; gpt-4o-mini). \todo{add o3 and one open-weights model}

\paragraph{Contamination caveat.} The LLM rows use a closed model that has
plausibly seen the underlying text-mined corpus during training; until a
contamination probe is run, LLM and agent rows should be read as
upper bounds and are excluded from our core claims.

\section{Results}
\label{sec:results}

\subsection{Recipe transferability of retrieval}

\begin{table}[t]
  \centering
  \caption{Transferability benchmark: formula-SRO@5 (primary metric,
  mean [95\% CI]) with element-SRO@5 for continuity, method consistency,
  and latency; $n=\nRetrievalQueries$ queries, corpus \corpusSize\
  materials. Every retriever is bounded below by the random floor and
  above by the oracle ceiling.}
  \label{tab:retrieval}
  \begin{tabular}{lcccc}
\toprule
Retriever & formula-SRO@5 & element-SRO@5 & MCR@5 & ms/query \\
\midrule
MAGPIE (SKY) & 0.162 [0.145, 0.180] & 0.325 & 0.502 & 34 \\
Random & 0.017 [0.014, 0.020] & 0.136 & 0.494 & 0 \\
Element Jaccard & 0.362 [0.335, 0.392] & 0.534 & 0.568 & 12 \\
Stoich Vector & 0.324 [0.301, 0.350] & 0.485 & 0.610 & 247 \\
TF-IDF & 0.197 [0.181, 0.214] & 0.363 & 0.578 & 59 \\
Hybrid (MAGPIE+Stoich RRF) & 0.268 [0.249, 0.289] & 0.413 & 0.564 & 318 \\
\midrule
\emph{Oracle ceiling} & 0.595 [0.575, 0.617] & -- & -- & -- \\
\bottomrule
\end{tabular}

\end{table}

Three observations from Table~\ref{tab:retrieval}. First, on the
non-circular formula-level metric, \textbf{raw element-set Jaccard
(\jaccFormulaSro) beats learned Magpie embeddings (\magpieFormulaSro) by
more than a factor of two}; the gap is larger than under the element-level
metric (\jaccElementSro\ vs \magpieElementSro), so the conclusion is not
an artifact of metric circularity. Property-oriented composition
embeddings distribute attention across dimensions irrelevant to synthesis.
Second, \textbf{every practical retriever is far from the ceiling}: the
oracle achieves \oracleFormulaSro, leaving the best baseline a mean regret
of \jaccMeanRegret\ --- and \jaccFracLargeRegret\ of queries lose more
than 0.2 formula-SRO to the ceiling. Recipe-aware retrieval is not a
solved subproblem but a wide-open one. Third, \textbf{the headroom is
concentrated}: for sol-gel queries the baseline--oracle gap is
\gapSolGel\ (\jaccSolGelFormulaSro\ vs \oracleSolGelFormulaSro) versus
\gapSolidState\ for solid-state --- invisible to prior solid-state-only
evaluations. \todo{add route-oracle ablation once reimplemented:
route-conditioning closes only a small fraction; remainder is
precursor-form selection}

\subsection{Downstream planning}

\begin{table}[t]
  \centering
  \caption{Planning benchmark on \nPlanningCases\ held-out targets.
  Precursor metrics as mean [95\% CI]. Cov.\ = fraction of cases with a
  scoreable temperature prediction.}
  \label{tab:planning}
  \begin{tabular}{lccccc}
\toprule
Method & Precursor Jaccard & Precursor F1 & Method acc. & Temp.\ MAE (\si{\celsius}) & Fail \\
\midrule
Rule baseline & 0.821 [0.787, 0.856] & 0.444 [0.360, 0.537] & 0.50 & 320 (84) & 0 \\
Retrieval (element Jaccard) & 0.754 [0.716, 0.788] & 0.448 [0.374, 0.522] & 0.59 & 320 (67) & 0 \\
Retrieval (Magpie) & 0.370 [0.303, 0.440] & 0.233 [0.174, 0.299] & 0.29 & 316 (45) & 42 \\
Direct LLM & 0.811 [0.777, 0.844] & 0.433 [0.352, 0.517] & 0.51 & 310 (84) & 0 \\
SKY agent & 0.813 [0.770, 0.855] & 0.497 [0.411, 0.583] & 0.53 & 345 (80) & 2 \\
\bottomrule
\end{tabular}

\end{table}

Table~\ref{tab:planning} gives the corrected downstream picture. The SKY
agent attains the highest precursor formula F1 \emph{point estimate}
(\skyFOne\ \skyFOneCI), above retrieval-copy (\retrJaccFOne) and the
ungrounded LLM (\llmFOne), but \textbf{no pairwise F1 difference involving
the agent survives Holm correction} (vs retrieval-copy:
$p_{\mathrm{raw}}=\skyVsRetrJaccFOnePRaw$,
$p_{\mathrm{Holm}}=\skyVsRetrJaccFOnePHolm$; vs ungrounded LLM:
$p_{\mathrm{raw}}=\skyVsLlmFOnePRaw$,
$p_{\mathrm{Holm}}=\skyVsLlmFOnePHolm$). At $n=\nPlanningCases$ the
benchmark cannot distinguish grounded-agent from retrieval-copy precursor
selection; a larger test set is required \todo{n=400 rerun}.

On the lenient element-level metric the zero-cost rule baseline
(\ruleElemJacc) is indistinguishable from the agent (\skyElemJacc):
precursor \emph{element} identity is nearly determined by the target
composition. Method accuracy must be read against the balanced design:
the 50/50 stratification destroys the corpus prior ($\sim$86\%
solid-state), and no method beats the better of (prior, chance) ---
element-Jaccard retrieval-copy is nominally best (\retrJaccMethodAcc).
On temperature, on the common subset ($n=\tempCommonN$) the rule
baseline's per-family median is best (\ruleTempMaeCommon~\si{\celsius})
and the agent worst (\skyTempMaeCommon~\si{\celsius}); differences among
non-failing methods are not significant, and we make no cross-protocol
temperature comparisons (Section 2). Magpie-based retrieval-copy fails on
\retrMagpieFailures\ of \nPlanningCases\ targets, echoing its
transferability deficit.

Taken together with Table~\ref{tab:retrieval}, the downstream evidence is
consistent with the augmentation-benchmark consensus~\cite{iclr2026benchmark}:
\textbf{retrieval quality is the operative bottleneck; agentic machinery
neither rescues weak retrieval nor demonstrably improves on
retrieval-copy at this sample size.}

\section{Discussion and limitations}

\paragraph{What the benchmark changes.} Transferability, oracle regret,
and route-stratified failure profiles convert ``our retriever finds
similar materials'' into testable, comparable claims, and identify a
concrete learnable target: route-conditioned precursor-\emph{form}
selection (nitrate/alkoxide/citrate vs oxide/carbonate, hydrates,
chelators), which dominates the sol-gel gap.

\paragraph{Limitations.} (i)~The corpus covers solid-state and sol-gel
routes only, inheriting text-mining biases~\cite{kononova2019text};
(ii)~formula-level Jaccard penalises hydrate/anhydrous mismatches
identically for all methods but understates chemical interchangeability;
(iii)~planning results use one closed model at $n=\nPlanningCases$ with a
possible training-data contamination confound; (iv)~temperature ground
truth is regex-extracted and route-heterogeneous --- our MAEs are not
comparable to per-operation solid-state protocols~\cite{prein2025ranking};
(v)~no experimental validation is attempted; (vi)~strong learned baselines
(He et al., Retrieval-Retro) are not yet run inside this harness ---
until then we claim a measurement gap, not superiority over any learned
retriever.

\section{Conclusions}

Retrieval grounds essentially all current synthesis-planning systems, yet
its recipe transferability has never been measured. Measuring it reveals
a wide oracle gap, concentrated in identifiable chemistry, and downstream
experiments where better retrieval --- not more agentic machinery --- is
the binding constraint. The benchmark, oracle tooling, frozen test sets,
and cached predictions are released for the community to close the gap we
quantify.

\section*{Data and code availability}
Code, benchmarks, frozen test sets, cached LLM outputs, and this
manuscript are available at
\url{https://github.com/KingaMas/synthesis-agent}. All prose statistics
are generated from committed result files by \texttt{paper/make\_tables.py}.
\todo{archive a release on Zenodo and add DOI}

\section*{Author contributions}
\todo{fill in CRediT roles with Hyunsoo}

\section*{Conflicts of interest}
There are no conflicts to declare.

\bibliographystyle{unsrtnat}
\bibliography{references}

\end{document}
